% -----------------------------------------------------------------------
% CHAPITRE 1 - INTRODUCTION (a adapter au contexte de votre universite)
% -----------------------------------------------------------------------
% A inclure dans votre .tex principal avec \input{chapitre1_introduction.tex}.
% Ce chapitre pose la problematique, les objectifs et la demarche generale
% sans anticiper les resultats (qui sont detailles au chapitre 3).
% -----------------------------------------------------------------------

\chapter{Introduction}
\label{ch:introduction}

\section{Contextualisation : la dynamique des paysages du Boeny}

\paragraph{Un territoire sous forte pression environnementale.} La région Boeny,
au nord-ouest de Madagascar, concentre plusieurs des grands enjeux écologiques du
pays : un littoral soumis au marnage de l'estuaire de la Betsiboka (un des plus
grands estuaires de l'océan Indien), une plaine alluviale à vocation agricole et
une mosaïque de forêts et de savanes. Ce « carrefour écologique » cumule toutefois
les pressions : urbanisation littorale, front pionnier agricole, défrichements et
instabilité de la bande côtière sous l'effet des cycles de décrue et de marée.

\paragraph{Un suivi insuffisant des trajectoires individuelles.} Les approches
classiques de télédétection se limitent souvent à deux images prises à deux dates
(état initial / état final), ce qui gomme les allers-retours intermédiaires
(ex. une parcelle qui passe par le sol nu avant de revenir en culture). Or, c'est
précisément cette \emph{séquence d'états successifs} — et non le seul état final
— qui permet de distinguer une trajectoire stable d'une trajectoire dynamique.

\paragraph{L'opportunité des séries temporelles décennales.} Depuis 2016, la
famille de produits \emph{Dynamic World} (Google Earth Engine) fournit chaque
année une carte globale de l'occupation du sol à 10 m de résolution, avec une
nomenclature d'une dizaine de classes (eau, forêt, végétation non forestière,
végétation inondée, cultures, bâti, sol nu). La combinaison de la haute
résolution spatiale (10 m) et du pas annuel sur dix ans (2016--2025) ouvre une
fenêtre d'observation inédite sur les trajectoires de changement à une échelle
fine, jusqu'alors réservée aux gros pixels (300 m--1 km) ou aux séries courtes.

\section{Problématique et hypothèses}

\paragraph{Problématique.} Face aux pressions qui s'exercent sur le Boeny, cette
recherche se demande : \emph{dans quelle mesure les trajectoires annuelles
d'occupation du sol, observées à 10 m puis agrégées à une résolution de 250 m sur
la période 2016--2025, permettent-elles de caractériser et de cartographier les
dynamiques de changement — stables, cycliques ou de transition — du paysage
boenien ?}

\paragraph{Hypothèses de travail.} Trois hypothèses structurent la démarche :

\begin{enumerate}
  \item \textbf{Des trajectoires typiques émergent spontanément.} Le bilan des
        transitions réalisé par regroupement statistique (k-means) fait émerger
        un petit nombre de trajectoires types (quelques unités), mêlant zones
        stables et zones dynamiques, sans qu'elles soient imposées a priori.
  \item \textbf{Les changements se produisent principalement en bordure des
        milieux stables.} La dynamique se concentre sur des franges de
        transition (forêt--savane, estuaire, front agricole) plutôt qu'au cœur
        des grandes unités homogènes.
  \item \textbf{Des associations de transitions sont récurrentes et localisées.}
        Certains enchaînements d'un état à un autre (ex. eau $\leftrightarrow$
        sol nu) sont statistiquement sur-représentés et spatialement circonscrits,
        signalant des processus spécifiques (vasières, estuaires, agriculture).
\end{enumerate}

\section{Objectifs et questions de recherche}

\paragraph{Objectif général.} Développer un pipeline reproductible permettant de
qualifier et de quantifier les trajectoires de changement de l'occupation du sol
du Boeny sur 2016--2025, puis d'en tirer une typologie spatiale et des règles
d'association.

\paragraph{Objectifs spécifiques.}

\begin{enumerate}
  \item \textbf{Agréger} les données 10 m (2016--2025) en une résolution d'analyse
        de 250 m, en appliquant un masque polygone au secteur d'étude, afin
        d'obtenir des trajectoires annuelles par cellule (1e phase).
  \item \textbf{Diagnostiquer et corriger le bruit} résiduel des séries
        (artefacts de capteur, changement aberrant d'une année sur l'autre) par
        un lissage spatio-temporel, en comparant l'avant et l'après.
  \item \textbf{Regrouper} les trajectoires en un nombre restreint de classes
        homogènes (clustering K-means) dont le nombre est validé par critère
        d'inertie et de silhouette.
  \item \textbf{Extraire} des règles d'association sur les transitions
        (Apriori), triées par lift, pour identifier les dynamiques de
        co-occurrence les plus informatives.
  \item \textbf{Visualiser} spatialement les résultats (cartes statiques et
        interactives) pour étayer l'interprétation géographique.
\end{enumerate}

\paragraph{Questions de recherche associées.}
\begin{enumerate}
  \item Les trajectoires du Boeny se regroupent-elles en classes bien séparées,
        et en combien ?
  \item Quelle part du territoire est stable / dynamique, et où ?
  \item Quelles transitions sont les plus associées (lesquelles s'enchaînent) —
        et ce constat est-il stable d'une année à l'autre ?
\end{enumerate}

\section{Démarche générale du mémoire}

Le mémoire s'organise en trois temps qui correspondent aux trois chapitres
suivants :

\begin{enumerate}
  \item \textbf{Chapitre 2 — Données et méthodes} : présentation du secteur
        (Boeny), des données Dynamic World 10 m 2016--2025, et de la chaîne de
        traitement (agrégation 10 m $\rightarrow$ 250 m, lissage, clustering,
        règles). Tous les choix méthodologiques sont exposés et justifiés.
  \item \textbf{Chapitre 3 — Résultats} : la typologie en six clusters, leur
        profil et leur répartition spatiale, puis les dix règles d'association
        (triées par lift) et leur lecture géographique.
  \item \textbf{Conclusion générale} : synthèse des réponses aux questions de
        recherche, limites (résolution d'analyse, lissage, choix de K) et
        perspectives (validation par données de terrain, extension à d'autres
        régions).
\end{enumerate}

\section{Reproductibilité}

L'ensemble des scripts (Python/R), les valeurs de paramètres (graine fixe
\texttt{SEED = 1}, \texttt{random\_state = 1}, \texttt{nstart = 25}) et les
modalités d'exécution sont documentés dans le dépôt associé et le guide
d'exécution (\texttt{GUIDE\_EXECUTION.md}). L'objectif est de rendre chaque
figure, chaque tableau et chaque chiffre du mémoire \emph{reproductibles à
l'identique} — une exigence centrale de la démarche scientifique.
